\section{Threat Model}
\label{sec:threat}

\subsection{Setting and Attacker}

\label{sec:attacker}

We call the person or organization that owns or operates the environment and
delegates a task to the agent the \textbf{user}. The agent runtime and tool
implementations are trusted, although content returned by a tool may come from
an attacker-controlled source.

The attacker is a legitimate participant in the environment, such as a vendor,
repository contributor, buyer, or seller. The attacker can place content only
in a source normally controlled by that role, and the source appears on the
benign execution path for a legitimate reason. The attack is planted before
execution: the attacker cannot modify the user or system prompt, observe the
agent's reasoning, intercept tool calls, or alter other participants' data. The
attacker knows the environment's public structure and the decisions for which
its source may foreseeably be consulted, but not the exact task issued by the
user, internal records unavailable to its role, or other participants' private
information. The attacker's goal is \textbf{task tampering}: a harmful deviation in
a decision that tasks consulting its source foreseeably involve (the action the
agent selects or one of its structured or free-text arguments), not a task of
its own. Not knowing the user's task, the attacker cannot ensure that a given
run reaches this decision. Each case therefore pairs the attack with a benign
task that does, and we measure the attack's effect given that alignment.

\subsection{Prompt-Injection Boundary}

\label{sec:ipi-boundary}

We define the prompt-injection boundary by the authority a source has over the
targeted decision. That authority may follow from the source holder's role,
the surrounding context, and the user task; our benchmark makes it explicit
rather than leaving the agent to infer it. A case is indirect prompt injection
when:

\begingroup
\setlist[enumerate]{topsep=2pt,itemsep=0pt,parsep=0pt,partopsep=0pt}

\begin{enumerate}
\def\labelenumi{\arabic{enumi}.}
\item
  the attacker controls content that the agent encounters while performing
  the benign task, and constructs it using only information available to its
  role or explicit fabrication;
\item
  the user request fixes the targeted decision or designates an agent-visible
  record or procedure that resolves it, and the attack's intended outcome is
  inconsistent with that result; and
\item
  accepting the planted content would produce a predefined harmful outcome.
\end{enumerate}

\endgroup

Conditions 1 and 3 follow the standard IPI setting and attacker goal
\citep{greshake2023ipi, liu2024formalizing, debenedetti2024agentdojo}; condition 2
makes the source's authority explicit in each case.

For example, in an inflated-refund case the user asks the agent to resolve a
damaged-item dispute and designates the platform order record as the basis for the refund amount; the
buyer's dispute evidence claims a \$1,074.99 refund although the record
authorizes \$899.99. If the task instead asks only for the buyer's stated
reason for the refund, a false damage claim comes from the source the task
designates for that answer, and condition 2 does not hold.

Context may establish that a source exceeded its authority even when the user
did not state the boundary. Because such context may be unverifiable at
decision time \citep{abdelnabi2026contextual}, each case instead states the boundary
in the user request and designates the source that resolves the decision; cases
without it are outside the benchmark, not necessarily outside
prompt injection.
